\documentclass[11pt, a4paper]{article}

% --- Packages ---
\usepackage[utf8]{inputenc}
\usepackage[margin=1in]{geometry} % Sets 1-inch margins
\usepackage{amsmath, amssymb, amsthm} % For math symbols and environments
\usepackage{hyperref} % Adds clickable links to the Table of Contents
\usepackage{xcolor} % For custom colors
\usepackage{tcolorbox} % For beautiful note/definition boxes
\usepackage{graphicx} % For inserting images
\usepackage{booktabs} % For professional-looking tables

% --- Hyperref Setup ---
\hypersetup{
	colorlinks=true,
	linkcolor=blue,
	filecolor=magenta,      
	urlcolor=cyan,
	pdftitle={Course Notes},
}

% --- Custom Environments ---
% A nice blue box for definitions or key concepts
\newtcolorbox{conceptbox}[1]{
	colback=blue!5!white,
	colframe=blue!75!black,
	fonttitle=\bfseries,
	title=#1,
	arc=4mm % Rounded corners
}

% --- Document Information ---
\title{\textbf{Course Title: Name of the Subject}}
\author{Your Name}
\date{\today} % Automatically inserts current date

% --- Main Body ---
\begin{document}
	
	\maketitle
	\thispagestyle{empty} % Removes page number from the title page
	\newpage
	
	% --- Table of Contents ---
	\tableofcontents
	\newpage
	
	% --- Lecture 1 ---
	\section{Lecture 1: Introduction to the Course}
	\subsection{Overview and Objectives}
	This is where you can start typing your general lecture notes. Use paragraphs to separate ideas. You can emphasize text using \textbf{bold text} or \textit{italics}.
	
	\subsection{Core Concepts}
	Here is an example of a bulleted list for quick facts:
	\begin{itemize}
		\item First key takeaway from the lecture.
		\item Second important point you need to remember.
		\item A sub-point looks like this:
		\begin{itemize}
			\item Detail regarding the second point.
		\end{itemize}
	\end{itemize}
	
	\begin{conceptbox}{Important Definition}
		This is a custom callout box. Use it to highlight formulas, definitions, or critical exam topics so they stand out when you review.
	\end{conceptbox}
	
	
	% --- Lecture 2 ---
	\section{Lecture 2: Advanced Topics}
	\subsection{Mathematical Formulations}
	If your course involves equations, LaTeX handles them beautifully. Here is an inline equation: $E = mc^2$. 
	
	If you want a standalone, numbered equation, use the equation environment:
	\begin{equation}
		f(x) = \int_{-\infty}^{\infty} \hat{f}(\xi)\,e^{2 \pi i \xi x}\,d\xi
	\end{equation}
	
	\subsection{Summary Table}
	You can also organize data using tables:
	\begin{center}
		\begin{tabular}{lrc}
			\toprule
			\textbf{Topic} & \textbf{Difficulty} & \textbf{Reviewed?} \\
			\midrule
			Lecture 1 & Easy & Yes \\
			Lecture 2 & Medium & No \\
			Lecture 3 & Hard & No \\
			\bottomrule
		\end{tabular}
	\end{center}
	
	
	% --- Appendix / Resources ---
	\newpage
	\section{Appendix: Additional Resources}
	Any extra links, reading lists, or formula sheets can go right here at the end of your document.
	
\end{document}